% -----------------------------------------------------------
\section{Offend, Offense, Offender, Offensive}
% -----------------------------------------------------------
	\label{OFFENSE}
	
% % ----------------------
% \subsection{See also}
% ----------------------

% ----------------------
\subsection{1828 American Dictionary of the English Language} % *1828
% ----------------------
	\label{OFFENSE-1828}
	
	% -----
	\subsubsection{Offend (transitive)}
	% -----

		\begin{compactitem}
			\item[1] To attack; to assail. [\textit{Not used.}]
			\item[2] To displease; to make angry; to affront. It expresses rather less than make angry, and without any modifying word, it is nearly synonymous with displease. We are \textit{offended} by rudeness, incivility and harsh language. Children \textit{offend} their parents by disobedience, and parents \textit{offend} their children by unreasonable austerity or restraint.
			\item[3] To shock; to wound; as, to \textit{offend} the conscience.
			\item[4] To pain; to annoy; to injure; as, a strong light \textit{offends} weak eyes.
			\item[5] To transgress; to violate; as, to \textit{offend} the laws. But we generally use the intransitive verb in this sense, with against; to \textit{offend} against the law.
			\item[6] To disturb, annoy, or cause to fall or stumble.
			\item[7] To draw to evil, or hinder in obedience; to cause to sin or neglect duty.
		\end{compactitem}

% % ----------------------
% \subsection{Oxford English Dictionary} % *OED
% % ----------------------
	% \label{OFFENSE-OED}

	% \begin{compactitem}
		% \item[]
	% \end{compactitem}

% ----------------------
\subsection{Scriptural Usage} % *SCRIPTURES
% ----------------------
	\label{OFFENSE-SCRIPTURES}

	% % -----
	% \subsubsection{Old Testament} % *OT
	% % -----
		% \label{OFFENSE-OT}
	
		% % --
		% \paragraph{KJV}
		% % --
			% \label{OFFENSE-OT-KJV}
	
			% \begin{tabular}{ll}
				% Total occurrences in the OT & 
			% \end{tabular}
			
			% \begin{compactdesc}
				% \item[]
			% \end{compactdesc}
			
		% % --
		% \paragraph{Hebrew Bible}
		% % --
			% \label{OFFENSE-HEBREW}
		
			% %
			% \subparagraph{\texorpdfstring{\he{sample word}}{}}
			% %
				% \label{PAGA} % whatever is in step bible without periods
			
				% \begin{compactdesc}
					% \item[HALOT , PDF ] \textbf{}
						% \begin{compactitem}
							% \item[1]
						% \end{compactitem}
					% \item[BDB , PDF ] \textbf{}
						% \begin{compactitem}
							% \item[1]
						% \end{compactitem}
					% \item[DCH PDF ] \textbf{}
						% \begin{compactitem}
							% \item[1]
						% \end{compactitem}
				% \end{compactdesc}
				
				% \begin{compactdesc}
					% \item[Some OT uses] \textbf{}
						% \begin{compactdesc}
							% \item[]
						% \end{compactdesc}
				% \end{compactdesc}
			
		% % --
		% \paragraph{Septuagint}
		% % --
			% \label{OFFENSE-LXX}
		
			% %
			% \subparagraph{\texorpdfstring{\gr{sample word}}{}}
			% %
		
	% -----
	\subsubsection{New Testament} % *NT
	% -----
		\label{OFFENSE-NT}
	
		% --
		\paragraph{KJV}
		% --
			\label{OFFENSE-NT-KJV}
			
			%
			\subparagraph{Offend}
			%	
	
				\begin{tabular}{ll}
					Total occurrences in the NT & 35
				\end{tabular}
				
				\begin{compactdesc}
					\item[{\hyperref[MATTHEW-11-6]{Matthew 11:6}}]
					\item[{\hyperref[MATTHEW-18-6]{Matthew 18:6--9}}]
				\end{compactdesc}
			
		% --
		\paragraph{Greek New Testament} % *GREEK
		% --
			\label{OFFENSE-GREEK}
		
			%
			\subparagraph{\texorpdfstring{\gr{σκανδαλίζω}}{}}
			%
				\label{SKANDALIZO}
			
				\begin{compactdesc}
					\item[BDAG 823a] \textbf{}
						\begin{compactitem}
							\item[1] \textbf{to cause to be brought to a downfall, \textit{cause to sin}} (the sin may consist in a breach of the moral law, or in the acceptance of false teachings)
							\item[1a] \gr{τινα} \textit{someone}...Passive, \textit{be led into sin}...The absolute passive can also mean \textit{let oneself be led into sin, fall away}
							\item[1b] Passive with \gr{ἔν τινι}, \textit{be led into sin, be repelled by someone} [said] of Jesus; by refusing to believe in him or by becoming apostate from him a person falls into sin \textbf{Mt 11:6; 13:57; 26:31, 33}
							\item[2] \textbf{to shock through word or action, \textit{give offense to, anger, shock}}
						\end{compactitem}
					\item[CGL 1268, PDF 1308] \textbf{}
						\begin{compactitem}
							\item[1] \textbf{cause} (with accusative someone) \textbf{to offend} or \textbf{sin}; passive, be led into sin
							\item[2] give offence to, \textbf{offend} ---someone; passive, take offence
						\end{compactitem}
					\item[LSJ , PDF ] \textbf{}
						\begin{compactitem}
							\item 
						\end{compactitem}
				\end{compactdesc}
				
				% \begin{compactdesc}
					% \item[Some NT uses] \textbf{}
						% \begin{compactdesc}
							% \item[]
						% \end{compactdesc}
				% \end{compactdesc}
		
	% % -----
	% \subsubsection{Book of Mormon} % *BOM
	% % -----
		% \label{OFFENSE-BOM}
	
		% \begin{tabular}{ll}
			% Total occurrences in the BOM & 
		% \end{tabular}
		
		% \begin{compactdesc}
			% \item[]
		% \end{compactdesc}
		
	% -----
	\subsubsection{Doctrine and Covenants} % *DC
	% -----
		\label{OFFENSE-DC}
	
		% \begin{tabular}{ll}
			% Total occurrences in the DC & 
		% \end{tabular}
		
		\begin{compactdesc}
			\item[{\hyperref[DC59-7]{DC 59:7}}]
			\item[{\hyperref[DC59-21]{DC 59:21}}]
		\end{compactdesc}
		
	% % -----
	% \subsubsection{Pearl of Great Price} % *PGP
	% % -----
		% \label{OFFENSE-PGP}
	
		% \begin{tabular}{ll}
			% Total occurrences in the PGP & 
		% \end{tabular}
		
		% \begin{compactdesc}
			% \item[]
		% \end{compactdesc}
		
% % ----------------------
% \subsection{Key Phrases} % *KEYPHRASES *PHRASES
% % ----------------------
	% \label{OFFENSE-PHRASES}

	% \begin{compactdesc}
		% \item[] \textbf{\#times} | list
			% % \begin{compactdesc}
				% % \item[Bible anchor verses] \phantom{.}
					% % \begin{compactdesc}
						% % \item[Romans 5:5] \gr{}
					% % \end{compactdesc}
				% % \item[Commentary] ---
			% % \end{compactdesc}
	% \end{compactdesc}
	
% % ----------------------
% \subsection{Bible Dictionaries} % *DICTIONARIES
% % ----------------------
	% \label{OFFENSE-BD}

% % ----------------------
% \subsection{Restoration Usage} % *RESTORATION
% % ----------------------
	% \label{OFFENSE-RESTORATION}

	% % -----
	% \subsubsection{Joseph Smith} % *JS
	% % -----
		% \label{OFFENSE-JS}
	
		% \begin{compactdesc}
			% \item[{\hyperref[KEY]{Date/Source}}]
		% \end{compactdesc}
	
	% % -----
	% \subsubsection{Brigham Young} % *BY
	% % -----
		% \label{OFFENSE-BY}
	
		% \begin{compactdesc}
			% \item[{\hyperref[KEY]{Date/Source}}]
		% \end{compactdesc}
	
% % ----------------------
% \subsection{Commentary and Studies} % *COMMENTARY *STUDIES
% % ----------------------
	% \label{OFFENSE-COMMENTARY}

	% % -----
	% \subsubsection{Key Points}
	% % -----
		% \label{OFFENSE-KEYS}
	
		% \begin{compactitem}
			% \item
		% \end{compactitem}
		
	% % -----
	% \subsubsection{Sample Study}
	% % -----
		% \label{OFFENSE-study}